\documentclass[first=firstname,last=lastname,company=comp,location=Dresden]{baarticle}
\usepackage{xcolor}
\usepackage{ifthen}

% Assertion macro
% see https://tex.stackexchange.com/questions/42328/testing-framework-api-for-latex
\newcommand\assert[2][]{%
    \ttfamily
    \def\result{#2}
    \space Test\,\stepcounter{tst} \thetst{} (Assert \detokenize{#1}):
    \ifthenelse{\equal{\test}{\result}}{%
        \textcolor{green}{Passed}%
    }{%
        \textcolor{red}{Failed}\\Expected: "#2"\\Got: "\test"%
    }
    \par%
}
\newcommand\assertUndef[1]{%
    \ttfamily
    \space Test\,\stepcounter{tst} \thetst{} (Assert \detokenize{#1}is not defined):
    \ifundef{#1}{%
        \textcolor{green}{Passed}%
    }{%
        \textcolor{red}{Failed}\\Expected: \detokenize{#1} to be undefined\\Got: \detokenize{#1} is "#1"%
    }
    \par%
}

% Test counter
\newcounter{tst}
\setcounter{tst}{0}

\makeatletter

\begin{document}
\mktitle{
    title=Test Report,
    img=example-image-a,
    course=science,
    number=1234,
    corrector={corrector 1,corrector 2},
    themedate=\today,
    returndate=\today,
    type=thesis
}
\clearpage
\edef\test{\thepage}\assert[title page increases page counter]{II}
\rmfamily
\mkblocknotice{
    signature=example-image-a,
    date=\today,
    location=Dresden
}
\mkaffirmation{
    signature=example-image-a,
    date=\today
}
\clearpage
\subsection*{Lengths}
\edef\test{\strip@pt\parindent}\assert[parindent is 0]{0}
\edef\test{\headrulewidth}\assert[headrulewidth is 0.4pt]{0.4pt}

\subsection*{Table of Contents\, List of Figures and Tables}
\edef\test{\cftsecdotsep}\assert[ToC lines filled with dots]{\cftdot}
\edef\test{\cftfigpresnum}\assert[LoF prefix is Abbildung]{Abbildung }
\edef\test{\cftfigaftersnum}\assert[LoF suffix is :]{:}
\edef\test{\cfttabpresnum}\assert[LoT prefix is Tabelle]{Tabelle }
\edef\test{\cfttabaftersnum}\assert[LoT suffix is :]{:}
\edef\test{\cftcodepresnum}\assert[LoC suffix is Quelltext]{Quelltext\space}
\edef\test{\cftcodeaftersnum}\assert[LoC suffix is :]{:}

\subsection*{Environments}

% Test environment captions
\csdef{@baenvtitle}{Title}
\csdef{@baenvsource}{}
\edef\test{\@bacaption}\assert[caption is Title (Quelle: eigene Darstellung)]{Title (Quelle: eigene Darstellung)}
\csdef{@baenvsource}{Source}
\edef\test{\@bacaption}\assert[caption is Title (Quelle: Source)]{Title (Quelle: Source)}
\csdef{@baenvsource}{}
\csundef{@baenvsource}
\csundef{@baenvtitle}

% We test every environment at least twice, to ensure no robust commands are used
\begin{bafigure}{Test1}
    Content2
\end{bafigure}
\begin{bafigure}[source=Source2,label=thesecond]{Test2}
    Content
\end{bafigure}
\begin{bafigure}[placement=b]{Test3}
    Content
\end{bafigure}
\assertUndef{\@baenvtitle}
\assertUndef{\@baenvsource}
\assertUndef{\@baenvlabel}
\assertUndef{\@baenvplacement}

% Latex insert \null after printing the number within \ref
% \edef\test{\ref{Test1}}\assert[label is @baenvtitle if not set explicitly]{1\null}
% \edef\test{\ref{thesecond}}\assert[label can be set]{2\null}

% We test every environment at least twice, to ensure no robust commands are used
\begin{batable}{Table1}
    The Table
\end{batable}
\begin{batable}[source=blup,placement=b]{Table2}
    The Table2
\end{batable}
% \edef\test{\ref{Table1}}\assert[label is @baenvtitle if not set explicitly]{1\null}
\assertUndef{\@baenvtitle}
\assertUndef{\@baenvsource}
\assertUndef{\@baenvlabel}
\assertUndef{\@baenvplacement}

\end{document}
